\label{chap:structure}

\section{Chains of CSZ factors}

\begin{lemma}[Chains]
  \label{lem:chain}
  \lean{CantorSzemeredi.MPS.CSZ_sSup_chain}
  \uses{def:csz, def:factor}
  Let $\{\Dbb_\alpha\}_{\alpha \in I}$ be a non-empty family of CSZ factors of $(X,\Abb,m,T)$
  totally ordered by inclusion, and let $\Dbb$ be the $\sigma$-algebra generated by their union.
  Then $\Dbb$ is a CSZ factor.
\end{lemma}

\begin{proof}
  \uses{def:csz, def:factor}
  $\Dbb$ is $T$-invariant since each $\Dbb_\alpha$ is.

  \emph{Approximation.}  For $E \in \Dbb$ and $\eta > 0$ there is $\alpha$ with
  $\norm{\one_E - \Exp(\one_E \mid \Dbb_\alpha)}_2 < \eta$.  Indeed $\bigcup_\alpha \Dbb_\alpha$ is
  an algebra of sets (any two members lie in a common $\Dbb_\alpha$), and
  $\bigcup_\alpha L^2(\Dbb_\alpha)$ is a linear subspace of $L^2(X)$ with closure $\mathcal H$.
  The sets $B \in \Dbb$ with $\one_B \in \mathcal H$ contain that algebra and form a monotone
  class, so by the monotone class theorem they are all of $\Dbb$.  Then
  $\Exp(\one_E \mid \Dbb_\alpha)$, being the orthogonal projection, is at least as close as any
  $g' \in L^2(\Dbb_\alpha)$.

  \emph{A set in a single factor.}  Fix $\Cantor$, $t \geq 1$, $i \geq 0$ and $E \in \Dbb$ with
  $m(E) > 0$.  Put $\kappa \coloneqq 1/(2(t+1))$, choose $\alpha$ with
  $g \coloneqq \Exp(\one_E \mid \Dbb_\alpha)$ satisfying $\norm{\one_E - g}_2^2 < \kappa^2 m(E)$,
  and set $H \coloneqq \{g > 1-\kappa\} \in \Dbb_\alpha$.  Then $m(H) > 0$: otherwise
  $1 - g \geq \kappa$ a.e.\ on $E$ and $\norm{\one_E-g}_2^2 \geq \kappa^2 m(E)$.

  \emph{Comparison, for every step.}  Let $d \in \ZZ$ and
  $x \in H_d \coloneqq H \cap T^{-d}H \cap \cdots \cap T^{-td}H$.  As $\Dbb_\alpha$ is
  $T$-invariant, $\Exp(T^{sd}\one_E \mid \Dbb_\alpha)(x) = g(T^{sd}x) > 1-\kappa$ for
  $0 \leq s \leq t$.  Since indicators take values in $\{0,1\}$,
  $\prod_{s} T^{sd}\one_E \geq 1 - \sum_s (1 - T^{sd}\one_E)$ pointwise, so on $H_d$,
  $\Exp(\prod_s T^{sd}\one_E \mid \Dbb_\alpha) > 1 - (t+1)\kappa = \tfrac12$.  As
  $H_d \in \Dbb_\alpha$ and the product is non-negative,
  \begin{equation}
    \label{eq:chain-comparison}
    m\bigl( E \cap T^{-d}E \cap \cdots \cap T^{-td}E \bigr) \geq \tfrac12 m(H_d) , \qquad d \in \ZZ .
  \end{equation}

  \emph{Conclusion.}  Take $d = b^i k_n$ in \eqref{eq:chain-comparison} and average over
  $M \leq n < N$: the liminf for $E$ is at least $\tfrac12$ the liminf for $H$, which is positive
  since $H \in \Dbb_\alpha$, $m(H) > 0$ and $\Dbb_\alpha$ is CSZ.
\end{proof}

\begin{remark}
  \eqref{eq:chain-comparison} is the \emph{only} point at which the CSZ property enters, and it
  holds for every step $d$ separately.  That is why the classical argument applies verbatim to
  the CSZ property with its uniform averages, its quantification over all Cantor sets, and its
  dilations: each of those is a statement about \emph{which steps $d$ are averaged}, and the
  comparison is uniform in $d$.  In Lean, factor \eqref{eq:chain-comparison} out as a standalone
  lemma about a single $d$; the rest is then formal.
\end{remark}

\section{The dichotomy}

\begin{lemma}[Dichotomy]
  \label{lem:dichotomy}
  \lean{CantorSzemeredi.MPS.weaklyMixingExtension_or_compactExtension}
  \notready
  \uses{lem:wiener}
  \uses{def:factor, def:cwm, def:compact-extension}
  Let $(X,\Abb,m,T)$ be an ergodic system satisfying the standing assumption, and let
  $\Dbb \neq \Abb$ be a factor.  Then either $\Abb$ is a weakly mixing extension of $\Dbb$, or
  there is a factor $\Ebb$ with $\Dbb \subsetneq \Ebb \subset \Abb$ which is a compact extension
  of $\Dbb$.
\end{lemma}

\begin{proof}
  Due to Furstenberg; see Moreira's notes, Proposition 8.6.1, for a proof in the present
  language.  The proof is for ergodic systems on compact metric spaces, which is the setting
  here, and uses the disintegration of $m$ over $\Dbb$.
\end{proof}

\begin{remark}[The one hard import]
  Lemma~\ref{lem:dichotomy} is quoted in the paper and is by far the largest single item in the
  formalization that is not written out there.  Two points about our definitions make it apply
  as stated, and both must be preserved in Lean:
  \begin{itemize}
    \item Moreira calls $f$ conditionally weakly mixing if $\abs{\cip{T^n f}{g}} \to 0$ in
      uniform Ces\`aro mean, in $L^2(\Dbb)$, for \emph{every} $g$ in the Hilbert module.  Our
      Definition~\ref{def:cwm} is the case $g=f$ with ordinary Ces\`aro means, hence
      \emph{implied} by his.  So a weakly mixing extension in his sense is one in ours, and the
      dichotomy as he proves it gives the dichotomy as we state it.
    \item Our Definition~\ref{def:compact-extension} of conditional compactness, with orbits
      indexed by $n \in \ZZ$ and approximants in the Hilbert module, is exactly his.
  \end{itemize}
  Recommended strategy: state Lemma~\ref{lem:dichotomy} as an axiom-free \texttt{theorem} with
  \texttt{sorry}, behind a clearly documented interface, and complete the rest of the development
  against it.  Discharging it is its own project.

  \smallskip
  \emph{What is in place towards it.}  The proof begins by taking $f$ in the module with
  $\Exp(f\mid\Dbb)=0$ which is \emph{not} conditionally weakly mixing, and producing an atom of
  the spectral measure $\sigma_G$ of $G = f \otimes f$.  The measure-theoretic half of that step
  is formalized: \texttt{exists\_atom\_of\_not\_tendsto\_cesaro\_zero} turns a non-vanishing
  Ces\`aro mean of $\abs{\widehat{\sigma}(n)} ^2$ into an atom.  Notably this needs only
  $(\sigma\times\sigma)(\Delta) = \int \sigma(\{\alpha\})\,d\sigma$, which is Fubini, and not the
  full identity $\sum_\theta \sigma(\{\theta\}) ^2 = (\sigma\times\sigma)(\Delta)$: an integral of
  a non-negative function is zero as soon as the function vanishes identically, so positivity of
  the diagonal exhibits an atom directly.  Together with \texttt{integral\_tensorSq\_comp}, which identifies $\widehat{\sigma_G}(n)$ with
  $\psi(n)$, and the converse estimate
  $\avg{N}\psi \to 0 \Rightarrow$ conditional weak mixing
  (\texttt{condWeaklyMixing\_of\_tendsto\_cesaro\_psi}, the triangle inequality
  \texttt{eLpNorm\_sum\_le} followed by Cauchy--Schwarz), \textbf{that first step is now
  proved}, as \texttt{exists\_atom\_of\_not\_condWeaklyMixing}.

  \smallskip
  \emph{The Hilbert--Schmidt step.}  Its first half is also formalized:
  \texttt{exists\_eigenvector\_of\_atom} turns an atom of the spectral measure at $\theta$ into a
  genuine eigenvector, $Uw = e(\theta)w$ with $w \neq 0$.  Von Neumann's mean ergodic theorem
  (which Mathlib has) makes the Birkhoff averages of the contraction $e(-\theta)U$ converge to
  the orthogonal projection onto its fixed points, and
  \texttt{tendsto\_norm\_birkhoff\_sq} --- the quadratic form identity \texttt{norm\_sum\_sq}
  followed by dominated convergence against the kernel $\abs{\mathrm{ces}_N(\alpha-\theta)} ^2$
  --- shows the limit has squared norm $\sigma(\{\theta\}) \neq 0$.

  The missing theory has since been supplied.  Mathlib has no Hilbert--Schmidt operators ---
  only Gram--Schmidt orthogonalization, which is unrelated --- so
  \texttt{CantorSzemeredi/Analysis/HilbertSchmidt.lean} develops them from scratch, written so
  that it can be upstreamed: the operator $T_k f(x) = \int k(x,y) f(y)\,d\nu(y)$ attached to a
  kernel $k \in L^2(\mu\otimes\nu)$, the bound $\norm{T_k f}_2 \le \norm{k}_2 \norm{f}_2$
  (\texttt{eLpNorm\_integralKernel\_le}), its packaging as a continuous linear map
  (\texttt{integralKernelL}), and the theorem that these operators are compact
  (\texttt{isCompactOperator\_integralKernelL}).  The compactness proof has no analogue in
  Mathlib either: a tensor product $g \otimes h$ gives a rank-one operator, and rank-one
  operators are compact (\texttt{isCompactOperator\_rankOne}, also absent upstream, which has
  the closure properties but no base case); the tensor products are total in
  $L^2(\mu\otimes\nu)$ because a function orthogonal to every one of them integrates to zero
  over every measurable rectangle, hence over every measurable set by Dynkin's lemma; and
  $k \mapsto T_k$ is continuous by the Hilbert--Schmidt bound, so compactness passes to the
  closure.

  \smallskip
  \emph{The fibrewise reading.}  This too is now formalized, in
  \texttt{CantorSzemeredi/Dynamics/Fibre.lean}.  Since
  $\lambda = \int m_x \otimes m_x \,dm(x)$ is an \emph{average} of product measures, Fubini for
  the joining (\texttt{lintegral\_joining}) makes the fibre integrals
  $\int \abs{H}^2 \,d(m_x\otimes m_x)$ average to $\norm{H}_{L^2(\lambda)}^2 < \infty$, so they
  are finite for $m$-a.e.\ $x$: the eigenfunction is, on almost every fibre, a Hilbert--Schmidt
  kernel (\texttt{ae\_memLp\_fibre}).  Each $m_x$ is a probability measure, so
  \texttt{isCompactOperator\_integralKernelL} applies fibre by fibre and gives
  \texttt{ae\_isCompactOperator\_fibre}: for a.e.\ $x$, $H$ defines a compact operator on
  $L^2(m_x)$.

  One point of care, recorded in that file: $H \in L^2(\lambda)$ gives only
  $\lambda$-a.e.\ strong measurability, which says nothing about any individual fibre --- a
  $\lambda$-null set can meet every fibre.  The lemmas therefore take a strongly measurable
  representative, which costs nothing, since \texttt{integralKernelL\_congr} says the operator
  depends only on the a.e.\ class of its kernel.

  \smallskip
  \emph{From fibre operators to the module.}  The fibre operators assemble into a single
  operator on the Hilbert module,
  $(\Phi_H a)(x) = \int H(x,y) a(y) \,dm_x(y)$, and
  \texttt{CantorSzemeredi/Dynamics/CondHS.lean} proves the estimate that governs it.  Two
  ingredients: \texttt{condKernel\_comp\_measure}, the disintegration $m = \int m_x \,dm(x)$
  (Mathlib proves this against the \emph{trimmed} measure, and since the kernel is
  $\Dbb$-measurable, \texttt{lintegral\_trim} removes the trim), whose consequence
  \texttt{ae\_ae\_condKernel} is what lets an a.e.\ statement about the base be used on almost
  every fibre; and \texttt{abs\_condIntegralOp\_le}, Cauchy--Schwarz on a single fibre.
  Together they give
  \[
    \norm{\Phi_H a}_{\Ebb|\Dbb} \;\leq\; C \, \norm{H}_{\mathrm{HS}, \text{fibre}}
    \qquad\text{whenever } \norm{a}_{\Ebb|\Dbb} \leq C
  \]
  (\texttt{condNorm\_condIntegralOp\_le}), and hence
  \texttt{condCompact\_of\_kernel\_net}: if the orbit of $f$ is realised as $\Phi_{K_n} a$ for a
  family of kernels admitting finite nets in the fibre Hilbert--Schmidt norm, then $f$ is
  conditionally compact in the sense of Definition~\ref{def:compact-extension}.  Applied to a
  \emph{difference} of kernels, the estimate is exactly what converts an approximation of
  kernels into an approximation of functions --- which is the shape \texttt{CondCompact} asks
  for, since its finite net must consist of module elements and must serve every $n \in \ZZ$
  at once.

  Note that these estimates carry their integrability hypotheses explicitly, collected in
  \texttt{IsFibreHS}.  That is not bureaucracy: $H \in L^2(\lambda)$ gives only a.e.\
  statements about the base, and a $\lambda$-null set can meet every fibre, so nothing
  fibrewise follows until it is pushed through \texttt{ae\_ae\_condKernel}.

  \smallskip
  \emph{The finite net.}  \texttt{condCompact\_of\_kernel\_net} is a reduction: it assumes a
  finite net of kernels.  Producing that net needs no measurable selection, contrary to what
  one might expect from the fibrewise compactness --- the orbit is not a general precompact
  set but a \emph{circle}.  Since $H$ is an eigenfunction, $H \circ R^n = e(n\theta) H$, so
  the orbit of kernels traces $(\cos n\theta, \sin n\theta)$ through the real span of
  $\Re H$ and $\Im H$.  The net is then total boundedness of a bounded subset of $\RR^2$,
  a fact about $\RR^2$ and not about the fibres, carried over by the two estimates
  \[
    \norm{c H}_{\mathrm{HS},\text{fibre}} = \abs{c}\,\norm{H}_{\mathrm{HS},\text{fibre}},
    \qquad
    \norm{H+K}_{\mathrm{HS},\text{fibre}}
      \leq \norm{H}_{\mathrm{HS},\text{fibre}} + \norm{K}_{\mathrm{HS},\text{fibre}} .
  \]
  The second is Minkowski applied twice --- inside each fibre for the rows, then across the
  fibre for the row norms (\texttt{hsFibre\_add\_le}); both uses go through
  \texttt{CantorSzemeredi/Analysis/SqrtIntegral.lean}, which states Cauchy--Schwarz and
  Minkowski for $\sqrt{\int f^2}$ rather than for \texttt{eLpNorm}, since a conditional norm
  is a real number depending on a base point which is then integrated again.  The outcome is
  \texttt{exists\_kernel\_net}, and combining it with the reduction gives
  \texttt{condCompact\_of\_param\_orbit}: an orbit of kernels drawn from a bounded
  two-parameter family, acting on a single module element, consists of conditionally compact
  functions.

  \smallskip
  \emph{The regularity is not an extra assumption.}  The measurability obligations follow
  from joint measurability of $H$: both $x \mapsto \norm{H(x,\cdot)}_{L^2(m_x)}$ and
  $\Phi_H a$ are Bochner integrals against the kernel, and Mathlib's
  \texttt{StronglyMeasurable.integral\_kernel\_prod\_right'} makes those measurable in the
  base point (\texttt{measurable\_hsRow}, \texttt{measurable\_condIntegralOp}).  The
  remaining per-fibre conditions reduce to conditions on the base: by the disintegration, a
  function in $L^2(m)$ is in $L^2(m_x)$ for a.e.\ $x$ (\texttt{ae\_memLp\_condKernel}).  So
  \texttt{IsFibreReg} and \texttt{IsFibreHS} are built by
  \texttt{isFibreReg\_of\_memLp\_hsRow} and \texttt{isFibreHS\_of\_memLp} from joint
  measurability, a condition on the rows, and global $L^2(m)$ bounds --- nothing fibrewise
  has to be checked by hand.

  \smallskip
  \emph{Constancy along fibres.}  The row condition --- that $\int H(z,y)^2\,dm_z(y)$ is
  finite for a.e.\ $z$ --- does \emph{not} follow from \texttt{ae\_memLp\_fibre}, which
  controls $\int\!\!\int H^2\,d(m_x \otimes m_x)$, an integral over a fibre \emph{pair}.
  Relating the two is the a.e.\ constancy of $z \mapsto m_z$ along fibres, and that is now
  proved (\texttt{ae\_ae\_integral\_condKernel\_eq}).  The argument needs neither countable
  generation nor measurable selection.  Fix a bounded measurable $g$ and put
  $G(x) = \int g\,dm_x$.  Since the kernel is $\Dbb$-measurable, so is $G$, hence $G$ is its
  own conditional expectation and $\int G\,dm_x = G(x)$ a.e.; the same applies to $G^2$,
  because that is $\Dbb$-measurable too.  Subtracting, the variance of $G$ over the fibre
  through $x$ vanishes, so $G$ is constant there.

  Applied one indicator at a time, this identifies two measures on $X \times X$: the graph
  measure $\rho = \int \delta_z \otimes m_z \,dm(z)$ equals the joining $\lambda$
  (\texttt{compProd\_condKernel\_eq\_joining}) --- drawing $z$ and then $y$ from $m_z$ has
  the same law as drawing a fibre and then both coordinates independently from it.  Only the
  functional form of constancy is used, one bounded $g = \one_B$ at a time, which is why no
  countable generating $\pi$-system is needed.  The row condition
  (\texttt{ae\_memLp\_row}) and the bound on the row norms (\texttt{memLp\_hsRow}) both drop
  out, so \texttt{isFibreReg\_of\_memLp\_joining} builds the whole of \texttt{IsFibreReg}
  from nothing but joint measurability of $H$ and $H \in L^2(\lambda)$.

  \smallskip
  \emph{The orbit identity.}  The conditional operator intertwines the dynamics:
  \[
    (\Phi_H a) \circ T^n \;=\; \Phi_{H \circ R^n}\,(a \circ T^n), \qquad R = T \times T
  \]
  (\texttt{condIntegralOp\_comp\_T} and \texttt{condIntegralOp\_comp\_iterate}).  The proof is
  short given what precedes it: both sides integrate row by row against the conditional
  measures, so by \texttt{compProd\_condKernel\_eq\_joining} both are integrals against
  $\lambda$, and the identity becomes the $R$-invariance of $\lambda$.  In particular no
  equivariance statement relating $m_{Tx}$ to $m_x$ is needed --- which matters, because the
  set-level \texttt{condKernel\_comp} does not give one without countable generation.

  \smallskip
  \emph{Reduction to finite-rank kernels.}  The orbit identity makes one thing explicit that
  was not visible before: along the orbit the transformation acts on the kernel \emph{and} on
  the function.  The kernel part is harmless --- for an eigenfunction $H \circ R^n$ runs
  around a circle, which \texttt{exists\_kernel\_net} handles --- but $a \circ T^n$ moves too,
  so \texttt{condCompact\_of\_param\_orbit}, which assumes a \emph{fixed} $a$, does not apply
  to an eigenfunction orbit; its statement now says so.

  What is needed instead is that $\Phi_H$ carries the bounded set $\{a \circ T^n\}$ to a
  conditionally totally bounded set, and \texttt{condCompact\_of\_kernel\_approx} reduces this
  to the finite-rank case: the conditional Hilbert--Schmidt bound transports an
  essential-sup approximation of \emph{kernels} to the \emph{functions}, the triangle
  inequality \texttt{condNorm\_add\_le} adds the two errors, and a finite-rank kernel acts by
  $\Phi_{\sum g_i \otimes h_i} b = \sum_i \cip{h_i}{b}_{\Dbb} \, g_i$
  (\texttt{condIntegralOp\_sum\_tensor}), with $\Dbb$-measurable coefficients.

  \smallskip
  \emph{What is still open, and a question about the source.}  The reduction leaves two
  hypotheses, and neither is a formality.

  \emph{(i)} It asks for $\norm{H-K}_{\mathrm{HS},\text{fibre}} \leq \delta$ \emph{a.e.\ on
  the base} with $K$ of finite rank, whereas $L^2(\lambda)$ density of tensor products gives
  smallness of the $L^2$ norm \emph{over} the base.

  \emph{(ii)} It asks for a finite net covering the finite-rank images $\sum_i c_i^n g_i$,
  where $c_i^n = \cip{h_i}{a \circ T^n}_{\Dbb}$ are bounded $\Dbb$-measurable \emph{functions}.
  Under the \emph{old} form of Definition~\ref{def:compact-extension} this meant a finite net
  for $\{c^n\}$ in the $L^\infty(\Dbb)$ ball, which does not exist.  With the definition
  corrected to Moreira's it dissolves: the index may depend on the fibre, and rounding
  $c^n(x)$ to a fixed finite grid pointwise is exactly such a choice, the finitely many grid
  points giving the finitely many approximants $\sum_i v_i g_i$.  See
  Remark~\ref{rem:compact-mismatch}.

  \smallskip
  \emph{(i) is settled by a cutoff.}  The bridge is the exact identity
  \[
    \int_X \norm{H}_{\mathrm{HS},\text{fibre}}^2 \,dm \;=\; \norm{H}_{L^2(\lambda)}^2
  \]
  (\texttt{lintegral\_hsFibre\_sq}), which is the disintegration applied twice --- once to
  unfold the fibre, once to recognise the row integrals as an integral against $\lambda$.  So
  $L^2(\lambda)$ approximation \emph{does} control $\norm{H-K}_{\mathrm{HS},\text{fibre}}$,
  but only in $L^2$ of the base; Chebyshev converts this to smallness outside a set of small
  measure.  The point is that the set may be taken in $\Dbb$, since
  $\norm{\cdot}_{\mathrm{HS},\text{fibre}}$ is $\Dbb$-measurable
  (\texttt{stronglyMeasurable\_hsFibre}) --- the kernel $x \mapsto m_x$ being
  $\Dbb$-measurable --- and a $\Dbb$-set is a union of fibres, so cutting the rows down to it
  simply multiplies the fibre norm by its indicator
  (\texttt{ae\_ae\_indicator\_eq}, \texttt{hsFibre\_indicator}).  The cut kernel then
  satisfies the bound \emph{everywhere}: this is \texttt{exists\_cutoff\_hsFibre\_le}, with
  $m(B^c) \leq \norm{H-K}_{L^2(\lambda)}^2/\delta^2$.  The discarded set is where the
  density in \emph{compact extension} absorbs the error, exactly as anticipated.

  \smallskip
  \emph{The density.}  The input to the cutoff --- that finite-rank kernels are dense in
  $L^2(\lambda)$ --- is now in place.  The repository already had
  \texttt{topologicalClosure\_span\_prodKernels}, but for complex kernels and a \emph{product}
  measure, and $\lambda$ is not a product.  The Dynkin argument behind it never used the
  product structure, only that measurable rectangles form a $\pi$-system generating the
  product $\sigma$-algebra, so
  \texttt{ae\_eq\_zero\_of\_forall\_setIntegral\_rectangle\_eq\_zero} is now stated for an
  arbitrary finite measure on $X \times X$ and any Banach target.  What does \emph{not}
  generalise is \texttt{memLp\_prod\_mul}: off a product measure
  $\norm{g \otimes h}_{L^2}$ is not controlled by $\norm{g}_2\norm{h}_2$, so the totality is
  proved for \emph{indicators} of rectangles --- bounded, hence in $L^2$ of any finite
  measure --- which is all that is needed
  (\texttt{topologicalClosure\_span\_rectKernels}).  \texttt{exists\_rectSum\_eLpNorm\_sub\_lt}
  extracts from this an explicit finite sum $\sum_i c_i \one_{A_i} \otimes \one_{B_i}$.

  Combining with the cutoff gives \texttt{exists\_finiteRank\_cutoff}: for every
  $\delta, \eta > 0$ a finite-rank kernel $K$ and a $\Dbb$-set $B$ with $m(B^c) \leq \eta$ on
  which $\norm{H-K}_{\mathrm{HS},\text{fibre}} \leq \delta$ everywhere.  This is exactly the
  hypothesis of \texttt{condCompact\_of\_kernel\_approx}, with the exceptional set $B^c$ left
  for the density in \emph{compact extension} to absorb.

  \smallskip
  \emph{What is still open.}  Two things, both now assembly rather than new ideas.  First,
  running \texttt{condCompact\_of\_kernel\_approx} against
  \texttt{exists\_finiteRank\_cutoff} and \texttt{condIntegralOp\_comp\_iterate}: the finite
  net for the finite-rank images, which is the grid-rounding described above, and the passage
  from the cut kernel back to $H$ through the density in \emph{compact extension}.  Second,
  the factor $\Ebb$ must be assembled from the conditionally compact functions.
\end{remark}

\section{The structure theorem for ergodic systems}

\begin{theorem}
  \label{thm:ecsz}
  \lean{CantorSzemeredi.MPS.CSZ_of_ergodic}
  \uses{lem:chain, lem:dichotomy, lem:extensions, lem:trivial-csz, def:csz}
  Every ergodic system satisfying the standing assumption is CSZ.
\end{theorem}

\begin{proof}
  \uses{lem:chain, lem:dichotomy, lem:extensions, lem:trivial-csz}
  Let $\mathbf K$ be the set of CSZ factors, partially ordered by inclusion.  It is non-empty
  (Lemma~\ref{lem:trivial-csz}) and every non-empty chain has an upper bound in $\mathbf K$,
  namely the $\sigma$-algebra generated by its union (Lemma~\ref{lem:chain}).  By Zorn's lemma
  $\mathbf K$ has a maximal element $\Dbb$.

  Suppose $\Dbb \neq \Abb$ and apply Lemma~\ref{lem:dichotomy}.  If $\Abb$ is a weakly mixing
  extension of $\Dbb$ then $\Abb$ is CSZ by Lemma~\ref{lem:extensions}(W).  Otherwise there is a
  compact extension $\Ebb \supsetneq \Dbb$, CSZ by Lemma~\ref{lem:extensions}(K).  Either way
  $\mathbf K$ contains a factor strictly larger than $\Dbb$, contradicting maximality.  Hence
  $\Dbb = \Abb$.
\end{proof}

\section{Reduction to ergodic systems on compact metric spaces}

\begin{proposition}
  \label{prop:reduction}
  \lean{CantorSzemeredi.MPS.CSZ_of_forall_ergodic_CSZ}
  \uses{thm:ecsz, def:csz}
  Suppose every ergodic system satisfying the standing assumption is CSZ.  Then every invertible
  measure preserving system is CSZ.
\end{proposition}

\begin{proof}
  \uses{thm:ecsz, def:csz}
  Fix $(X,\Abb,m,T)$, a Cantor set $\Cantor$ of base $b$, integers $t \geq 1$, $i \geq 0$, and
  $A \in \Abb$ with $m(A) > 0$.

  \emph{A compact metric model.}  Let $\Omega \coloneqq \{0,1\}^{\ZZ}$ with the product topology and
  shift $(\tau\omega)(n) = \omega(n+1)$, a homeomorphism.  Define $\Phi : X \to \Omega$ by
  $\Phi(x)(n) \coloneqq \one_A(T^n x)$; each coordinate is measurable and
  $\Phi \circ T = \tau \circ \Phi$, so $\mu \coloneqq \Phi_\ast m$ is $\tau$-invariant.  With
  $\widetilde A \coloneqq \{\omega : \omega(0)=1\}$ we have
  $\Phi^{-1}(\tau^{-d}\widetilde A) = T^{-d}A$ for every $d \in \ZZ$, hence
  \begin{equation}
    \label{eq:model}
    m\Bigl( \bigcap_{s=0}^{t} T^{-d_s}A \Bigr) = \mu\Bigl( \bigcap_{s=0}^{t} \tau^{-d_s}\widetilde A \Bigr) ,
    \qquad d_0,\dots,d_t \in \ZZ ,
  \end{equation}
  and $\mu(\widetilde A) = m(A) > 0$.  So it suffices to prove the positive liminf for
  $\widetilde A$ in $(\Omega,\mathcal B,\mu,\tau)$.

  \emph{Ergodic decomposition.}  For $\mu$-a.e.\ $y$ there is an ergodic $\tau$-invariant Borel
  probability $\mu_y$ with $y \mapsto \mu_y(B)$ measurable and $\mu(B) = \int \mu_y(B)\,d\mu(y)$.
  Write $c_n(\nu) \coloneqq \nu(\bigcap_{s=0}^{t}\tau^{-s b^i k_n}\widetilde A)$, so
  $c_n(\mu) = \int c_n(\mu_y)\,d\mu(y)$.  Since $0 < \mu(\widetilde A) = \int \mu_y(\widetilde A)d\mu(y)$,
  the set $Y_+ \coloneqq \{y : \mu_y(\widetilde A) > 0\}$ has positive measure, and for $y \in Y_+$
  the system $(\Omega,\mathcal B,\mu_y,\tau)$ is ergodic, invertible and on a compact metric
  space, hence CSZ: $c(y) \coloneqq \unifliminf \unifavg c_n(\mu_y) > 0$.

  \emph{Passing the liminf through the integral.}  Let
  $c \coloneqq \unifliminf \unifavg c_n(\mu)$ and choose intervals $[M_k,N_k)$ with
  $N_k - M_k \to \infty$ along which $\avg{M_k \leq n < N_k} c_n(\mu) \to c$.  The function
  $\ell(y) \coloneqq \liminf_k \avg{M_k \leq n < N_k} c_n(\mu_y)$ is measurable and
  $\ell(y) \geq c(y)$, since a liminf along one sequence of intervals of growing length is at
  least the liminf over all such intervals.  By Fatou (the integrands being non-negative),
  $c \geq \int \ell \, d\mu \geq \int_{Y_+} \ell \, d\mu > 0$.
\end{proof}

\begin{remark}
  Fatou's lemma is the only place where the uniformity of the CSZ property costs anything: a
  uniform liminf is not a liminf along a single sequence, so Fatou cannot be applied to it
  directly.  The device above --- extract one sequence of intervals realizing the uniform
  liminf, apply Fatou along it, and compare with the uniform liminf fibrewise --- is what the
  Lean proof must reproduce.
\end{remark}

\section{The correspondence principle and the main theorems}

\begin{theorem}[Furstenberg correspondence principle]
  \label{thm:furstenberg-correspondence}
  \lean{CantorSzemeredi.furstenberg_correspondence}
  \notready
  \uses{def:banach-density, def:mps}
  Let $A \subset \ZZ$ with $\dstar(A) > 0$.  Then there are an invertible measure preserving
  system $(X,\Abb,m,T)$ and $\tilde A \in \Abb$ with $m(\tilde A) = \dstar(A)$ such that for every
  finite collection of integers $r_1,\dots,r_j$,
  \begin{equation}
    \label{eq:correspondence}
    \dstar\bigl( A \cap (A-r_1) \cap \cdots \cap (A-r_j) \bigr)
      \geq m\bigl( \tilde A \cap T^{-r_1}\tilde A \cap \cdots \cap T^{-r_j}\tilde A \bigr) .
  \end{equation}
\end{theorem}

\begin{proof}
  Furstenberg, Theorem 1.1.  The symbolic model of Proposition~\ref{prop:reduction} is half of
  the construction; the other half is a weak-$\ast$ limit of averages of point masses along
  intervals realizing $\dstar(A)$.  Formalizing it needs weak-$\ast$ compactness of the
  shift-invariant probability measures on $\{0,1\}^{\ZZ}$, which Mathlib supports.
\end{proof}

\begin{proof}[Proof of Theorem~\ref{thm:main-ergodic}]
  \proves{thm:main-ergodic}
  \uses{thm:ecsz, prop:reduction, lem:csz-functions, def:csz}
  By Proposition~\ref{prop:reduction} and Theorem~\ref{thm:ecsz}, every invertible measure
  preserving system is CSZ.  The case $i=0$ of \eqref{eq:csz}, together with
  Lemma~\ref{lem:csz-functions} for the passage from sets to non-negative functions, is exactly
  Theorem~\ref{thm:main-ergodic}.
\end{proof}


